\section{Planned CFD Validation Workflow for the Drag Discrepancy}
The initial BARAM cross-check established that the drag discrepancy for \texttt{IntendedValidation2} is pressure-dominated rather than primarily viscous. That finding narrows the next validation step, but it does not by itself isolate the underlying mechanism. A structured follow-on study is therefore required before a stronger aerodynamic conclusion can be defended.

\subsection{Validation Objectives}
The validation campaign should answer four specific questions. First, is the CFD drag level numerically converged with respect to mesh and wall treatment? Second, where on the aircraft does the excess drag originate? Third, is the excess pressure drag associated with blended-centerbody form effects, wing-body interference, or localized separation? Fourth, once the CFD case is tightly matched to the repository cruise condition, how much of the discrepancy remains attributable to the low-order model itself rather than to setup mismatch?

\subsection{Current Limitation of the Available Case}
The currently available BARAM case is useful as an initial cross-check, but it has one important limitation: the aircraft surface exists as a single wall patch. Consequently, the present solution can decompose drag into pressure and viscous parts, but it cannot directly attribute drag to the centerbody, outer wing, control surfaces, or body-wing junction. That limitation prevents direct validation of the centerbody and interference hypotheses from the existing case alone.

\subsection{Required Case Modifications}
The first required change is a patch-split geometry export. At minimum, the aircraft should be remeshed with separate named surfaces for the centerbody, outer wing, and control surfaces. A more informative split would isolate the inner wing and, if practical, the wing-body junction region as well. The second required change is exact condition matching between CFD and the repository cruise point, including geometry revision, reference area, density, velocity, angle of attack, symmetry convention, and moment reference point. The third required change is the addition of per-patch force and force-coefficient monitors so that pressure and viscous drag can be reported by region rather than only for the aircraft as a whole.

\subsection{Proposed Study Matrix}
The first case in the validation matrix should be a matched baseline rerun at the repository cruise condition using thesis-level reference quantities in the CFD force monitor. The second should be a three-level mesh-convergence study on that same matched case. The third should be a patch-resolved rerun using the split aircraft surface so that drag can be assigned to the centerbody, wing, and control-surface regions. The fourth should focus on flow visualization, including surface pressure coefficient, center-plane velocity slices, wake slices, and streamlines near the blended centerbody and wing root. The fifth should be a geometry-ablation study in which the centerbody is simplified or removed while holding the outer-wing condition as constant as practical. If the drag remains sensitive or ambiguous after those steps, a sixth study should compare at least one alternate turbulence model to the standard \(k\)-\(\epsilon\) result.

\subsection{Criteria for Interpreting the Results}
If the patch-resolved results show that most of the excess pressure drag is concentrated on the centerbody or wing-body junction, then the centerbody/interference hypothesis is supported. If the same regions also show adverse-pressure signatures, recirculation, or separated wake structures, then a separation-supported pressure-drag interpretation becomes defensible. If, however, the drag level shifts substantially with mesh refinement or turbulence-model choice, then the discrepancy should be reported as not yet physically isolated. In that case, the correct thesis claim is not that a single mechanism has been identified, but that the integrated workflow has successfully exposed a high-value validation problem and organized the evidence needed to resolve it.
